% !TEX root = ../main.tex

\chapter{Présentation fonctionnelle du logiciel FPG}
\label{chap:presentation}

Avant de détailler la mise en œuvre technique de l'historisation, il est
nécessaire de présenter le fonctionnement du logiciel FPG tel que
découvert et manipulé au cours du stage. Cette présentation permet de
comprendre le périmètre exact de la mission~: quelles interfaces sont
concernées, quelles informations doivent être tracées, et selon quelle
logique métier.

Le schéma ci-dessous offre une vue d'ensemble du parcours utilisateur au
sein du logiciel, selon le profil de l'utilisateur connecté~: il sera
détaillé étape par étape tout au long de ce chapitre.

\begin{figure}[H]
  \centering
  \imgwlink{schema_parcours_utilisateur.png}{\textwidth}
  \caption{Vue d'ensemble du parcours utilisateur}
  \label{fig:parcours-utilisateur}
\end{figure}

\section{Deux interfaces distinctes}

Le logiciel FPG repose sur deux interfaces distinctes, destinées à deux
profils d'utilisateurs différents.

\subsection{Interface Adhérent}

L'interface \textbf{Adhérent} est accessible aux utilisateurs publics,
membres d'une structure adhérente (une entreprise cotisante). Elle leur
permet~:

\begin{itemize}
  \item de rédiger le brouillon d'une demande de financement~;
  \item de joindre les pièces justificatives nécessaires~;
  \item de soumettre la demande, ce qui entraîne la création du dossier
        avec le statut \texttt{TRANSMIS}~;
  \item de consulter l'état d'avancement du dossier et d'échanger via une
        messagerie, selon les droits exposés par l'API.
\end{itemize}

\subsection{Interface Collaborateur}

L'interface \textbf{Collaborateur} (dite \og{}office\fg{}) est réservée
aux agents du FPG en charge du traitement des dossiers. Elle leur permet
notamment~:

\begin{itemize}
  \item d'instruire le dossier~;
  \item de modifier les informations relatives à la formation, aux
        stagiaires ou au montage financier~;
  \item de faire évoluer le statut du dossier~;
  \item de gérer les règlements, les documents administratifs
        (\textit{MAPC (Modificative d'attestation de prise en Charge)}), les contrôles, les notes internes et les
        affectations aux collaborateurs.
\end{itemize}

C'est exclusivement sur cette interface Collaborateur que la
fonctionnalité d'historisation a été implémentée, dans la mesure où
c'est elle qui concentre l'essentiel des actions de modification
apportées à une demande de financement.

\section{Architecture technique globale}

Le schéma ci-dessous présente l'architecture technique globale du
logiciel FPG, telle qu'observée et manipulée au cours du stage~: le
frontend Angular, l'API REST Django, le serveur d'authentification
Keycloak, la base de données PostgreSQL, le stockage de fichiers MinIO,
ainsi que le serveur de génération de rapports JasperReports.

\begin{figure}[H]
  \centering
  \imgw{Schéma_architecture_FPG.pdf}{\textwidth}
  \caption{Schéma d'architecture technique du logiciel FPG}
  \label{fig:architecture-fpg}
\end{figure}

Cette architecture constitue le socle sur lequel repose l'ensemble des
fonctionnalités présentées dans ce chapitre, et plus particulièrement la
fonctionnalité d'historisation détaillée au chapitre suivant.

\section{Connexion et redirection}

L'accès au logiciel se fait via une page de connexion unique. Selon les
identifiants saisis, l'utilisateur est automatiquement redirigé vers
l'interface qui lui correspond  -  Adhérent ou Collaborateur.

\begin{figure}[H]
  \centering
  \imgwlink{page_de_connexion.png}{0.7\textwidth}
  \caption{Page de connexion commune aux deux interfaces}
  \label{fig:page-connexion}
\end{figure}

\section{Liste des demandes de financement}

Une fois connecté à l'interface Collaborateur, l'utilisateur accède à la
liste des demandes de financement, point d'entrée principal de son
activité quotidienne.

\begin{figure}[H]
  \centering
  \imgwlink{onglet_demande_de_financement.png}{0.9\textwidth}
  \caption{Liste des demandes de financement, interface Collaborateur}
  \label{fig:liste-demandes}
\end{figure}

Cette liste présente, pour chaque dossier, les informations
principales~: le numéro de dossier, la date de réception, le type de
financement, les dispositifs complémentaires choisis, le millésime, le
numéro Tahiti et le numéro CPS de la structure, l'intitulé de l'action de
formation, l'organisme de formation prestataire, les dates de début et
de fin de formation, le collaborateur assigné au dossier, l'état interne,
ainsi que le statut du dossier. Un bouton d'action permet enfin d'accéder
aux différents onglets de la demande de financement, chacun d'entre eux
pouvant faire l'objet de modifications à historiser.

\section{Cycle de vie d'une demande de financement}

\subsection{Flux des statuts}

Une demande de financement traverse, au cours de son traitement, une
succession de statuts~: de sa création (\texttt{TRANSMIS}) jusqu'à son
règlement final, en passant par les étapes d'instruction, de préparation
de règlement, etc. Le schéma ci-dessous représente l'ensemble des
transitions possibles entre statuts.

\begin{figure}[H]
  \centering
  \imgwlink{schema_des_flux_status.png}{0.85\textwidth}
  \caption{Flux des statuts d'une demande de financement}
  \label{fig:flux-statuts}
\end{figure}

\subsection{Verrouillage des onglets selon le statut}

Le statut d'un dossier détermine quels onglets sont modifiables à un
instant donné. Certains onglets ne se déverrouillent, par exemple, que
lorsque le dossier atteint le statut \og{}Préparation règlement\fg{}. La
compréhension fine de cette matrice de verrouillage s'est révélée
essentielle pour la réalisation des tests de la fonctionnalité
d'historisation~: elle seule permet de savoir, pour un statut donné,
quels onglets sont effectivement testables, et donc de vérifier de
manière exhaustive que chaque modification possible est correctement
tracée.

\begin{figure}[H]
  \centering
  \imgwlink{verrou_selon_statut.png}{0.85\textwidth}
  \caption{Matrice de verrouillage des onglets selon le statut du dossier}
  \label{fig:verrou-statuts}
\end{figure}

Cette matrice a ainsi constitué un outil de travail quotidien tout au
long du développement et des campagnes de tests de la fonctionnalité
d'historisation, détaillée au chapitre suivant.